
\chapter{Платёжный шлюз изнутри}\label{ch:gateway}

Один запрос продавца~--- \texttt{POST} к \texttt{/v1/payment\_intents}
на оплату заказа. Снаружи это обычный JSON поверх HTTPS, внутри
тот же запрос проходит через последовательность решений,
и~ошибка в~любом звене меняет риск, конверсию или
PCI scope. Здесь интеграция превращается в~архитектуру.

\section{Анатомия шлюза}\label{sec:gw-anatomy}

Платёжный шлюз (payment gateway) превращает один запрос продавца
в~корректную, наблюдаемую авторизацию.

\definitionnote{Шлюз / процессор / оркестратор}{Три термина,
  которые часто используют как взаимозаменяемые, но на деле они
  описывают разные уровни абстракции. \textbf{Платёжный шлюз}
  (payment gateway) принимает платёжные запросы от продавца по
  высокоуровневому API и переводит их в~протоколы карточных сетей и
  банков-эквайеров. \textbf{Платёжный процессор} (payment processor)
  непосредственно обрабатывает авторизацию, клиринг и~расчёт и
  часто имеет прямое подключение к~карточным сетям. \textbf{Платёжный
  оркестратор} (payment orchestrator) маршрутизирует транзакции между
  несколькими шлюзами или процессорами по заданным правилам.
  Разделение между ролями часто размыто. Крупные международные провайдеры
  платёжных услуг (PSP) вроде Stripe или Adyen совмещают все три
  уровня в~одном продукте; на российском рынке аналогичную совмещённую
  роль выполняют ЮKassa, CloudPayments, T-Бизнес, ранее Тинькофф,
  Сбер для платформ, Best2Pay, RBK.money и~другие.}

Внутри платёжного шлюза выделяются модуль приёма API-запросов,
модуль оценки риска, модуль маршрутизации, хранилище токенов
и~модуль авторизации, отправляющий запрос банку-эквайеру
или напрямую карточной сети. Порядок их вызова не случаен~---
он задаёт и архитектурные свойства системы, и~её PCI scope.

Продавец отправляет HTTPS-запрос на API шлюза, например,
\texttt{POST} к \texttt{/v1/payment\_intents} либо
\texttt{/v1/charges} в~более старой версии API, с JSON-телом,
содержащим сумму, валюту, данные карты или токен и метаданные
заказа. Модуль приёма валидирует сообщение, проверяет
аутентификацию по API key или OAuth token и идемпотентность,
подробнее в~разделе~\ref{sec:gw-reliability}, затем передаёт
нормализованное внутреннее представление транзакции модулю оценки риска.

Модуль оценки риска применяет к~транзакции набор правил и~моделей:
проверки частоты повторов (velocity checks), сопоставление
геолокации, отпечаток устройства, а~в~зрелых системах ещё и
скоринговую модель вероятности мошенничества. На выходе~---
одно из трёх решений: пропустить транзакцию дальше, отправить
её на дополнительную проверку по 3-D Secure или отклонить.
Устройство антифрод-систем разобрано в~главе~\ref{ch:fraud};
здесь важен порядок. Проверка риска работает \emph{до} авторизации,
потому что подозрительную транзакцию дешевле остановить заранее,
чем потом разбирать чарджбэк.

Если транзакция прошла оценку риска, она попадает в~модуль
маршрутизации, который определяет, через какого эквайера
или процессор отправить авторизационный запрос.
Логика маршрутизации остаётся одной из самых коммерчески
значимых частей шлюза~--- §\ref{sec:gw-routing}.

Если продавец передал полный PAN, шлюз заменяет
его токеном из хранилища. PAN сохраняется в~зашифрованном
хранилище, а~дальше по внутреннему конвейеру и в~ответе продавцу
циркулирует только токен. Это сужает CDE (cardholder data environment)
до одного компонента и радикально сокращает PCI scope всей
системы~--- механизм подробно разобран в~главе~\ref{ch:pci-dss}.
Если продавец с~самого начала прислал
сетевой токен (DPAN), хранилище может не участвовать, но шлюз
всё равно проверяет ограничения домена токена и~при необходимости
запрашивает у TSP актуальную криптограмму для конкретной
транзакции~\cite{emvco-payment-tokenisation-guide}
(см. также главу~\ref{ch:tokenization}).

Модуль авторизации формирует сообщение ISO~8583 или его
эквивалент в API-протоколах новых эквайеров, отправляет запрос
по защищённому каналу банку-эквайеру, получает ответ, преобразует
код ответа в~формат API шлюза и возвращает результат продавцу.
Ожидается прежде всего ответ эмитента через карточную сеть.

\begin{figure}[tbp]
  \centering
  \includefiguresvg[width=.45\textwidth]{assets/figures/ch27-payment-gateway/gateway-components}
  \caption{Компоненты платёжного шлюза.}\label{fig:gateway-components}
\end{figure}

Главное свойство шлюза видно на этом минимальном маршруте:
каждый следующий компонент получает более узкое и~лучше
определённое представление транзакции, чем предыдущий.
Из этого вырастают и требования к~изоляции данных, и PCI scope,
и ограничения на отказные режимы.

\section{Интеллектуальная маршрутизация}\label{sec:gw-routing}

При нескольких эквайерах шлюз решает, куда
отправить авторизационный запрос, чтобы не проиграть
одновременно в~стоимости, доле одобрений и времени ответа.

\subsection*{Оси оптимизации}

Первая ось~--- \textbf{стоимость}. Тарифы у~эквайеров разные,
а interchange fee зависит от типа карты, MCC ТСП, региона
эмитента и~десятков других параметров.
Модуль маршрутизации выбирает эквайера с~минимальной
стоимостью для конкретной комбинации условий~--- например,
направляет внутрироссийскую транзакцию по карте Мир через
эквайера с~более выгодным тарифом НСПК, а трансграничную~---
через другого, у~которого ниже валютная наценка.

Вторая ось~--- \textbf{доля одобрений}. На одних и~тех же
BIN-диапазонах эквайеры показывают разную долю одобрений.
Причины бывают техническими или коммерческими: качество
подключения к~сети, скорость ответа, репутация ТСП
у~конкретного эквайера, локальный или трансграничный маршрут.
Накопив статистику по BIN, стране и~сумме, модуль маршрутизации
направляет транзакцию туда, где вероятность одобрения выше.
Для крупных объёмов даже небольшая разница в~качестве маршрута
превращается в~заметное изменение выручки.

Третья ось~--- \textbf{латентность}. Продавцу в~реальном времени
важно, чтобы ответ пришёл быстро. Если один эквайер стабильно
отвечает заметно быстрее другого, модуль маршрутизации при
прочих равных предпочтёт первый, потому что каждая лишняя секунда
в~сценарии оплаты снижает конверсию.

Эти оси неизбежно конфликтуют. Самый дешёвый эквайер
может иметь худшую долю одобрений, а~самый быстрый~--- оказаться
самым дорогим. Задача модуля маршрутизации~--- найти рабочий
оптимум в~ограничениях, заданных продавцом: \enquote{минимизировать
  стоимость при доле одобрений не ниже X\% и p95 латентности
не более Y~мс}. Продвинутые системы решают такую задачу
как адаптивное распределение трафика; в~литературе этот подход
описывают моделью multi-armed bandit.

\subsection*{Каскадная маршрутизация}

Каскадную маршрутизацию часто путают с~интеллектуальной,
хотя она решает другую задачу. Если первый маршрут вернул отказ,
допускающий ещё одну попытку по другому маршруту, шлюз
отправляет транзакцию второму эквайеру. Термин
\enquote{мягкий отказ} здесь нужно использовать аккуратно,
у~схем он может иметь узкий нормативный смысл. Mastercard, например,
требует мягкий отказ в~странах со строгой аутентификацией клиента
(strong customer authentication, SCA) только тогда, когда
CNP-операции нужна аутентификация, но её нет,~--- а~не как
общий синоним любого отказа~\cite{mc-tpr}. Логика каскада
остаётся прежней. Отказ одного маршрута ещё не означает,
что транзакция в~принципе невозможна. Другой эквайер может
получить иной ответ от того же эмитента, особенно когда
его авторизационный запрос отличается: другой BIN эквайера,
другой MCC, другое значение Terminal ID.

\importantnote{%
  Каскадная маршрутизация и обычный повтор решают разные задачи.
  Каскад отправляет транзакцию \emph{другому} эквайеру
  после отказа, допускающего такой манёвр; повтор идёт \emph{тому же}
  эквайеру, обычно спустя некоторое время.
  Путаница между этими сценариями и приводит к~дублям. Если
  шлюз повторил запрос тому же эквайеру, а~первый ответ лишь задержался,
  эмитент увидит два авторизационных запроса и~может одобрить
  оба. Каскад сам по себе не создаёт дубля на стороне первого
  эквайера, но создаёт другой риск: задержавшийся ответ от
  первого маршрута может прийти после того, как второй
  успел получить одобрение. Грамотная реализация должна
  немедленно отправлять отмену первой авторизации, если та
  материализовалась с~опозданием.

  Каскад работает только для части отказов, которые эквайер
  или схема считают потенциально устранимыми. К~явно неверным
  реквизитам~--- \texttt{14} Invalid Card Number или
  \texttt{54} Expired Card --- это не относится~\cite{mc-tpr}.
  Сами схемы ограничивают избыточные повторные попытки;
  конкретные лимиты зависят от сценария и~правил
  конкретной схемы~\cite{mc-tpr}. Модуль маршрутизации
  обязан различать эти два класса отказов по коду ответа
  и~включать каскад только в~первом случае.
}
\definitionnote{Каскад}{Повторная попытка через \emph{другого}
  эквайера или маршрут. Это не обычный повтор тому же партнёру
и~не универсальный ответ на любой отказ.}

Интеллектуальная маршрутизация отвечает на вопрос \enquote{куда лучше
вести первую попытку}, каскадная~--- \enquote{что делать после мягкого
отказа}. У~этих логик разные цели и~разные отказные режимы;
смешивать их опасно.

\section{Интеграционные модели}\label{sec:gw-integration}

Способ подключения продавца к~шлюзу определяет, как
делится ответственность за карточные данные и~каков
PCI scope продавца.

\definitionnote{Размещённая платёжная страница (HPP)}{Платёжная страница,
  размещённая на стороне шлюза. ТСП перенаправляет туда покупателя
и не принимает карточные данные в~свой периметр.}
\subsection*{Размещённая платёжная страница (перенаправление)}

В~модели HPP продавец инициирует платёжную сессию через API шлюза
и~получает URL, на который перенаправляет покупателя. Покупатель
вводит данные карты на странице шлюза, шлюз обрабатывает транзакцию
и возвращает покупателя на сайт продавца с~результатом. Параллельно
шлюз отправляет продавцу асинхронное уведомление с~финальным
статусом.

Главное преимущество HPP --- минимальный PCI scope. Карточные
данные не проходят через серверы продавца, поэтому
он обычно попадает в~модель SAQ A, а~не в~более широкие
сценарии, где элементы платёжной страницы частично исходят
от самого продавца~\cite{pci-faq-1291,pci-faq-1293}. Для небольших
ТСП без выделенной команды безопасности это часто
единственный реалистичный вариант.

Недостаток~--- потеря контроля над пользовательским сценарием.
Перенаправление прерывает поток оплаты, страница шлюза может
выбиваться из визуального языка сайта, обратный переход
после оплаты не всегда проходит гладко. Сайт продавца при такой
схеме не исчезает из PCI scope полностью: PCI SSC отдельно
уточняет, что redirect mechanism оставляет сайт продавца
в~области оценки~\cite{pci-faq-1332}.

\subsection*{Встроенная форма (iFrame / JavaScript SDK)}

Продавец встраивает в~свою страницу платёжную форму шлюза~---
iFrame или готовый JavaScript-компонент.
Визуально форма выглядит как часть сайта продавца, но технически
данные карты вводятся внутри iframe, контролируемого шлюзом,
и отправляются напрямую на серверы шлюза, минуя серверы продавца.

Модель сохраняет больше контроля над интерфейсом~--- продавец
стилизует форму, не разрывает сценарий оплаты и~претендует
на SAQ A или SAQ A-EP. Разница зависит
от деталей реализации и от того, проходят ли карточные данные
через JavaScript продавца хотя бы транзитом~\cite{pci-faq-1291,pci-faq-1293}.

\practicenote{%
  В PCI DSS v4.0.1 страницы оплаты находятся под отдельным
  контролем~--- требования 6.4.3 и 11.6.1 посвящены их защите.
  В~январе 2025 года PCI SSC обновил SAQ A, убрав эти два
  требования и заменив их упрощённым критерием~--- продавец
  подтверждает, что его сайт не подвержен атакам через скрипты.
  Обновлённые критерии разъяснены в FAQ 1588~\cite{pci-dss-4,pci-faq-1588}.
  Даже если карточные данные вводятся в iFrame шлюза, вредоносный
  скрипт на странице продавца может перехватить ввод через
  подложный оверлей или keylogger до того, как данные попадут
  в iFrame. При встроенной форме продавцу нужно доказуемо
  контролировать скрипты, влияющие
  на payment page~\cite{pci-faq-1588}.

}
\subsection*{Прямая интеграция через API (server-to-server)}

Продавец собирает данные карты на своей стороне через
собственную форму и отправляет их на API шлюза
в~серверном запросе. Это даёт полный контроль над
интерфейсом и произвольную предварительную логику:
проверку адреса, предавторизацию, разделение платежа
между несколькими получателями. Платой служит
существенно более широкий PCI scope~--- карточные данные проходят
через серверы продавца, и~сценарий уже не относится к~льготным
моделям SAQ A / A-EP~\cite{pci-faq-1291,pci-faq-1293,pci-saq}.

Этот путь выбирают крупные ТСП и платёжные платформы,
у~которых есть PCI DSS Level 1 и~для которых прямой
контроль над данными~--- конкурентное преимущество.
Сюда входит собственное защищённое хранилище токенов
и оптимизация логики повторов на уровне PAN.

\practicenote{%
  Выбор интеграционной модели~--- это и~бизнес-решение.
  Начинать разумно с iFrame или SDK~--- приемлемый пользовательский
  опыт при минимальном PCI scope. На server-to-server переходят,
  когда ограничения встроенной формы блокируют
  конкретный сценарий, например, Card-on-File с~прозрачной
  для покупателя токенизацией или каскадную маршрутизацию
  на уровне PAN.
}

\section{Надёжность и наблюдаемость}\label{sec:gw-reliability}

Надёжность шлюза определяется дисциплиной обращения
с~неопределённостью, когда ответ запаздывает, внешний вызов
теряется или продавец нажимает \enquote{повторить}.

\subsection*{Идемпотентность}

Повторный вызов с~тем же набором параметров обязан давать тот же
результат, что и~первый, без побочных эффектов.
Если продавец отправил запрос на авторизацию, не получил ответа
из-за таймаута или сетевого сбоя и~повторил его с~тем же idempotency key,
шлюз возвращает результат первой попытки, а~не создаёт
вторую транзакцию~\cite{stripe-idempotency}.

Для этого шлюз хранит устойчивое сопоставление между
\path{idempotency_key} и \path{transaction_id}
в~хранилище с~гарантированной консистентностью~--- обычно
в~реляционной базе с~уникальным индексом по ключу.
Срок жизни ключа должен перекрывать окно, в~котором
ещё возможны безопасные повторы одной и~той же операции,~---
от сетевого сбоя сразу после ответа до позднего повторного
запроса клиента.

\importantnote{%
  Идемпотентность на уровне шлюза не заменяет идемпотентность
  на уровне эквайера. Если шлюз отправил авторизацию,
  не получил ответа и~отправил повторно, эквайер может
  обработать оба запроса. Шлюз либо использует
  DE~11 (STAN) + DE~37 (RRN) для дедупликации
  на стороне эквайера, либо реализует сценарий
  \enquote{authorize-then-check}~--- при таймауте отправляет не повтор,
  а~запрос статуса у~эквайера, когда такой механизм предусмотрен
  протоколом, и~только при статусе \enquote{не найдено} повторяет
  авторизацию.

}
\subsection*{Обработка таймаута}

Таймаут опасен тем, что отсутствие ответа~--- это \emph{не}
отказ. Если шлюз отправил авторизационный запрос эквайеру
и~не получил ответа в~течение порога, возможны как минимум
три исхода: (а)~запрос не дошёл; (б)~запрос дошёл,
эмитент одобрил транзакцию, но ответ потерялся на обратном пути;
(в)~запрос дошёл, эмитент отказал, ответ тоже потерялся.
В~момент таймаута шлюз не различает эти случаи.

Наихудший исход~--- (б): деньги заблокированы на карте
держателя, эмитент считает транзакцию одобренной, продавец
не знает об этом и~может повторить платёж, создав двойное списание.
Отмена~--- штатный способ снять неопределённость в~карточном
обмене. Mastercard, например, отдельно требует \textit{automatic reversal},
если после \textit{authorization request} ответ не был получен~\cite{mc-tpr}.
После таймаута шлюз сначала закрывает неопределённое состояние
и затем решает, что возвращать продавцу.

\subsection*{Сбой шлюза между приёмом ответа и~записью}

Таймаут описывает ситуацию, когда ответ не пришёл. Обратная
проблема опаснее~--- ответ пришёл, эмитент одобрил транзакцию,
но шлюз упал до того, как персистировал результат. Состояние
транзакции расщепляется. Эмитент удерживает деньги на карте
держателя, а~шлюз после перезапуска видит запись со статусом
\texttt{sent\_to\_acquirer} и~не знает исхода.

Если шлюз просто повторит авторизацию, эмитент получит два
холд на одну сумму~--- двойное списание. Если шлюз вернёт
продавцу ошибку, деньги останутся заблокированными без
товара. Стандартное решение~--- двухфазная запись.
До отправки запроса эквайеру шлюз создаёт запись
\texttt{pending} с idempotency key и STAN. После перезапуска
фоновый процесс выбирает все \texttt{pending}-записи старше
порога 5--15~с и выполняет запрос статуса (\textit{status inquiry}) к~эквайеру по STAN
и RRN. Если эквайер подтверждает одобрение, шлюз переводит
запись в \texttt{approved} и~ставит вебхук в~очередь. Если
эквайер не знает о~транзакции, шлюз записывает
\texttt{failed} и освобождает продавца для повторной попытки.
Когда запрос статуса недоступен (не все эквайеры его
реализуют), шлюз отправляет отмену, гарантированно
снимая холд, и возвращает продавцу ошибку с~рекомендацией
повторить~\cite{mc-tpr}.

Без этого механизма каждый перезапуск шлюза порождает
зависшие холды: держатели карт видят их как
списания, продавцы~--- как потерянные заказы. На масштабе
тысяч транзакций в~минуту даже секундный сбой создаёт десятки
расхождений, разбор которых вручную через эквайера занимает дни.

\subsection*{Дедупликация}

Идемпотентность решает часть проблемы. Дедупликация
от двойной обработки одной и~той же транзакции работает
на нескольких уровнях. На уровне API она опирается
на idempotency key. На уровне внутреннего конвейера~---
на уникальность комбинации STAN + RRN + amount + timestamp,
которая отлавливает дубли, проскочившие мимо идемпотентности:
например, при каскадной маршрутизации, когда два эквайера
получают запрос с~разными STAN, но за одну и~ту же покупку.
На уровне денежных состояний дедупликацию берёт на себя
внутренний реестр обязательств из главы~\ref{ch:internal-ledger},
а~на уровне внешних файлов и банковской выписки работает
сверка из главы~\ref{ch:reconciliation}.

\subsection*{Вебхуки и~гарантия доставки}

Платёжный поток редко заканчивается в~момент синхронного
API-ответа. Результат транзакции возвращается продавцу
двумя путями: синхронно, в~ответе на запрос, и асинхронно,
через обратный HTTP-вызов (вебхук) на указанный URL.
Вебхук нужен потому, что синхронный ответ может
не дойти, а~после первого ответа появляются
новые внешние события: завершился 3-D Secure challenge,
пришло подтверждение, сработал возврат, поступило позднее
обновление статуса от процессора.

Публичные API современных PSP --- от международных Stripe и
Adyen до российских ЮKassa, CloudPayments, T-Бизнес и~Сбер
для бизнеса~--- строят доставку вебхуков по принципу
\enquote{как минимум один раз}. Эндпоинт должен быть доступен
по HTTPS, система ждёт ответ 2xx, при неуспехе повторяет
доставку, а~продавец проверяет подпись события, прежде чем
доверять содержимому запроса. У Stripe подпись
передаётся в HTTP-заголовке \texttt{Stripe-Signature: t=...,v1=...},
где \texttt{t} --- timestamp подписи, а \texttt{v1} ---
HMAC-SHA256 от конкатенации timestamp и тела запроса
с~использованием секрета подписи
вебхуков~\cite{stripe-webhooks}.\footnote{Правила подписи и~формат
вебхуков у~разных PSP различаются. У~одних подпись HMAC-SHA256
в HTTP-заголовке, у~других~--- JWT-токен или собственная схема;
список повторов и таймаутов тоже разный. Перед интеграцией
с~конкретным PSP сверяйтесь с~его документацией, а~не с~типовыми
\enquote{Stripe-подобными} ожиданиями.}

\practicenote{%
  Вебхуки у~шлюза приходят как минимум один раз, поэтому обработчик
  на стороне продавца обязан быть идемпотентным~--- до изменения состояния
  заказа нужно проверить, не обработан ли уже этот \path{event_id}.

}
\subsection*{Наблюдаемость}

При тысячах транзакций в~минуту метрики должны показать деградацию
раньше, чем она затронет значительную долю трафика.
Операционная команда непрерывно отслеживает несколько групп сигналов.

Доля одобрений~--- процент одобренных транзакций
от общего числа попыток, с~разбивкой по эквайеру, BIN-диапазону,
стране эмитента, MCC и~сумме. Падение доли одобрений
за короткий интервал~--- сигнал к~немедленному расследованию:
причиной бывает как сбой у~эквайера, так и~массовая атака
скомпрометированными картами.

Время ответа отслеживается через p50, p95 и p99 обработки
запроса, измеренного от получения HTTPS-запроса продавца
до отправки ответа. Неожиданный рост p99 указывает
на проблемы с HSM, медленный ответ эквайера или перегрузку
собственного модуля маршрутизации.

Доля системных ошибок включает транзакции,
завершившиеся таймаутом, сетевой или внутренней ошибкой сервера,
а~не бизнес-ответом вроде \texttt{approved} или \texttt{declined}. Когда эта
метрика становится заметной на фоне обычного трафика, команда
имеет дело с~полноценным инцидентом.

Доставку вебхуков продавцам показывает доля уведомлений,
дошедших с~первой попытки, выявляя либо проблемы
на стороне продавцов, либо перегрузку собственной очереди
доставки.

\section{От REST API к ISO~8583}\label{sec:gw-api-to-iso}

Пока разработчик смотрит на транзакцию как на JSON-объект
с~полями \texttt{amount}, \texttt{currency},
\path{payment_method} и \texttt{description}, банк-эквайер
и карточная сеть видят ту же операцию как сообщение
ISO~8583. В~публичных API современных PSP --- международных
Stripe и Adyen или российских ЮKassa, CloudPayments, T-Бизнес,
Сбер для бизнеса~--- это буквально JSON поверх HTTPS с~чётко
описанным поведением запроса и~ответа~\cite{stripe-api-reference,yookassa-api,cloudpayments-api}.\footnote{Названия
эндпойнтов и набор полей у~разных PSP свои. У Stripe это
\texttt{POST~/v1/payment\_intents} и \texttt{POST~/v1/charges},
у~ЮKassa --- \texttt{POST~/v3/payments}, у CloudPayments ---
\texttt{/payments/cards/charge} для одностадийного платежа или
\texttt{/payments/cards/auth} для двухстадийного, с~поддержкой
JSON и form-encoded.
Иллюстрация ниже намеренно использует обобщённый
\texttt{POST~/charges} как нейтральный пример, а~не дословный
эндпойнт конкретного PSP.} Шлюз транслирует одно
представление в~другое; о~самом стандарте см.
главу~\ref{ch:iso8583}.

Для российского ТСП есть ещё один путь, не имеющий аналога
в~карточном API уровня Stripe~--- оплата через СБП по динамическому
QR. На уровне продукта это тот же API-вызов вида
\texttt{POST~/payments} с~указанием способа оплаты СБП;
дальше PSP создаёт платёжную сессию у~банка-партнёра и QR-код
в~формате СБП, а~шлюз возвращает продавцу ссылку или payload
для отображения. Внутри маршрут идёт не в ISO~8583, а~в OpenAPI
ОПКЦ СБП (§\ref{sec:sbp-protocol}). Этот сценарий~---
повод проектировать платёжный шлюз не как \enquote{транслятор
в ISO~8583}, а~как абстрактный оркестратор методов оплаты,
для которого ISO~8583 --- один из исходящих протоколов.

Конвейер обработки одного \texttt{POST~/charges} внутри шлюза.

\begin{enumerate}
  \item \textbf{Приём и валидация}~--- HTTP-слой проверяет
    API-ключ, парсит JSON, валидирует обязательные поля
    (\texttt{amount}~>~0, \texttt{currency} по ISO~4217,
    наличие \path{payment_method}). Если передан
    \texttt{idempotency-key}, шлюз ищет кэшированный ответ
    и~при совпадении возвращает его без повторной обработки.
  \item \textbf{Скоринг фрода}~--- шлюз вычисляет фрод-скор сам или
    запрашивает у~внешнего сервиса. Если скор выше порога
    отказа, транзакция отклоняется до контакта с~эквайером,
    экономя стоимость авторизационного запроса
    (см. главу~\ref{ch:fraud}).
  \item \textbf{Маршрутизация}~--- модуль маршрутизации выбирает
    эквайера по BIN-диапазону, валюте, MCC и~текущей доле
    одобрений каждого канала. Если основной канал деградировал,
    каскадное правило переключает трафик на резервный маршрут.
  \item \textbf{Извлечение PAN}~--- если \path{payment_method}
    ссылается на сохранённый токен, шлюз обращается к~хранилищу
    vault за FPAN или DPAN. Vault возвращает карточные данные
    в~оперативную память процесса; на диск они не попадают
    (см. главу~\ref{ch:pci-dss}).
  \item \textbf{Сборка ISO~8583}~--- шлюз формирует MTI~0100,
    заполняет bitmap и data elements согласно IFS выбранного
    эквайера. Маппинг полей описан ниже.
  \item \textbf{Отправка эквайеру}~--- сериализованное сообщение
    уходит по TCP/TLS-соединению из пула. Шлюз запускает
    таймер ожидания ответа; верхний предел задают правила
    платёжной системы и согласованный с~эквайером SLA.
    Для НСПК внутрироссийский авторизационный таймаут
    исчисляется десятками секунд; конкретный пункт сверяется
    по действующей редакции Правил платёжной системы
    \enquote{Мир}~\cite{nspk-mir-rules}.
  \item \textbf{Парсинг ответа}~--- ответ MTI~0110 парсится,
    DE~39 (Response Code) транслируется в JSON-статус. Шлюз
    персистирует результат и обновляет запись транзакции.
  \item \textbf{Ledger event}~--- шлюз публикует событие
    \texttt{authorization.created} во внутреннюю шину. Модуль
    реестра обязательств фиксирует pending debit/credit
    (см. главу~\ref{ch:internal-ledger}).
  \item \textbf{Вебхук}~--- событие ставится в~очередь доставки
    продавцу, подписывается HMAC-ключом эндпоинта и отправляется
    на URL обратного вызова с~повтором при неуспехе.
  \item \textbf{Синхронный ответ}~--- HTTP-слой возвращает JSON
    продавцу с~полями \texttt{status}, \texttt{id},
    \path{decline_code} при отказе, timestamp.
\end{enumerate}

Конкретные значения p50/p95 латентности конвейера зависят
от множества факторов: расположения шлюза, эквайера и~эмитента,
протокола сети, нагрузки HSM, политики повторов. Публичных
агрегированных бенчмарков по российскому рынку нет, поэтому
числа в~книге не приводятся; ориентир для собственного
измерения~--- разделение полного времени ответа на сегменты:
внутренняя обработка шлюза, ожидание эквайера, ожидание
сети и~эмитента.

Сопоставление JSON-полей с data elements относится к~шагу~5.
Продавец отправляет запрос с~суммой \texttt{1500}, валютой
\texttt{RUB}, номером карты и~сроком действия. Шлюз транслирует
это в~сообщение MTI~0100 (Authorization Request), заполняя
десятки полей данных:

\begin{itemize}
  \item \texttt{amount} 1500 $\to$ DE~4 Amount, Transaction~---
    12-символьное числовое поле, выровненное нулями,
    \texttt{000000150000} в~минорных единицах, копейках;
  \item \texttt{currency} \texttt{RUB} $\to$ DE~49 Currency Code,
    Transaction~--- \texttt{643}, числовой код рубля по ISO 4217;
  \item \texttt{card.number} $\to$ DE~2 Primary Account Number~---
    PAN, до 19 цифр;
  \item \path{card.exp_month}, \path{card.exp_year} $\to$
    DE~14 Date, Expiration~--- \texttt{YYMM};
  \item внутренний ID терминала ТСП $\to$ DE~41 Card Acceptor
    Terminal Identification~--- 8-символьный код;
  \item MCC ТСП $\to$ DE~18 Merchant Type~--- 4-значный код.
\end{itemize}
Шлюз также генерирует поля, которых нет в API-запросе
продавца, но требуемые протоколом~--- DE~11 (STAN), уникальный
в~коротком окне, обычно сутки или смена, 6-значный
номер транзакции, DE~7 (Transmission Date
and Time), DE~12 (Time, Local Transaction), DE~22 (Point
of Service Entry Mode)~--- способ ввода данных карты
и~канал предъявления,~--- и DE~25 (Point of Service Condition Code).

Если продавец использует сетевой токен вместо PAN, шлюз заполняет
DE~2 токенизированным PAN (DPAN), а криптограмму от TSP
передаёт в EMV data или в~специфичном для схемы поле, зависящем
от конкретного протокола маршрута~\cite{emvco-payment-tokenisation-guide}
(см. также главу~\ref{ch:tokenization}).

Ответ эмитента проходит обратный путь. DE~39 (Response Code),
двухсимвольный код, шлюз транслирует в JSON-поле
\texttt{status} со значениями \texttt{succeeded}, \texttt{declined},
\path{requires_action} и дополняет человекочитаемым
описанием. Шлюз \emph{теряет} часть
информации: ISO~8583 различает десятки причин отказа
(\texttt{05} Do Not Honor, \texttt{51} Insufficient Funds,
\texttt{61} Exceeds Withdrawal Limit), тогда как API шлюза группирует их
в~несколько категорий. Чтобы отличить
\enquote{недостаточно средств} от \enquote{карта заблокирована},
разработчику нужно смотреть на поле \path{decline_code}
или его аналог в~дополнение к \texttt{status};
подробнее об интерпретации кодов ответа см.~\cite{mc-tpr}
и~главу~\ref{ch:decline-management}.

\practicenote{%
  Многие шлюзы предоставляют подробный протокольный лог транзакции,
  raw response, со значениями data elements DE 39,
  DE 44 (Additional Response Data), DE 55 (EMV data). Если шлюз
  этого не показывает, а~эквайер возвращает непонятный отказ,
  запросите у~эквайера trace log по STAN и RRN --- это единственный
  способ понять, что именно произошло на уровне протокола.
}

Трансляция охватывает авторизацию и другие операции. Возврат, void
и подтверждение не сводятся к~одному универсальному набору сообщений
для всех шлюзов и эквайеров, конкретное соответствие зависит
от стадии операции и протокола маршрута. На уровне API
это несколько продуктовых действий, а~на уровне карточного
обмена~--- последовательность authorisation, отмен, financial
messages и clearing updates, которые потом оказываются
во внутреннем реестре обязательств и в~сверке.
